% BHU PYQ -- Transcribed & Corrected
% Paper: GGRSE-21 / GGRSE21: Basics of Remote Sensing
% Exam: Under Graduate Semester II Examination, 2024-25 (NEP)
% Category: Skill Enhancement Course (SEC)
% Department: Geography

\documentclass[11pt,a4paper]{article}
\usepackage[a4paper,top=2.5cm,bottom=2.5cm,left=2.8cm,right=2.8cm,headheight=14pt]{geometry}
\usepackage{amsmath,amssymb}
\usepackage[T1]{fontenc}
\usepackage[utf8]{inputenc}
\usepackage{lmodern,microtype,fancyhdr,enumitem,hyperref,lastpage,parskip}

\hypersetup{
    colorlinks=true,
    urlcolor=black,
    linkcolor=black,
    pdftitle={GGRSE21: Basics of Remote Sensing -- BHU Sem II 2024-25 (NEP SEC)}
}

\pagestyle{fancy}
\fancyhf{}
\renewcommand{\headrulewidth}{0.4pt}
\renewcommand{\footrulewidth}{0.4pt}
\fancyhead[L]{\small   \textit{Banaras Hindu University}}
\fancyhead[R]{\small   \textit{GGRSE-21: Basics of Remote Sensing (SEC)}}
\fancyfoot[C]{\small Page  \thepage\ of \pageref{LastPage}}


\newlist{parts}{enumerate}{2}
\setlist[parts,1]{label=(\alph*),leftmargin=*,itemsep=4pt,topsep=2pt}

\newcommand{\pts}[1]{\hfill   \textit{[#1]}}


\begin{document}


\begin{center}
  {\large\bfseries BANARAS HINDU UNIVERSITY}\\[2pt]
  {
\normalsize Under Graduate --- Semester II Examination, 2024-25 (NEP)}\\[6pt]
  \rule{0.8\linewidth}{0.5pt}\\[6pt]
  {\large\bfseries GEOGRAPHY}\\[4pt]
  {\large\bfseries Skill Enhancement Course (SEC)}\\[4pt]
  {\large\bfseries Paper Code: GGRSE-21 : Basics of Remote Sensing}\\[4pt]
  {
\normalsize Time: 2:30 Hours\hfill Maximum Marks: 70}\\[2pt]
  {\small REF NO. CECONF/N/2024-25/024}
\end{center}

\rule{\linewidth}{0.4pt}

\smallskip



\noindent   \textit{Instructions: Answer   \textbf{five} questions in all including Question No.~1 which is compulsory. Answer each part of Question No.~1 in approximately 200 words and remaining questions in approximately 1000 words. Illustrate your answers with suitable maps and diagrams. All questions carry equal marks (14 marks each).}

\bigskip




\noindent   \textbf{Question 1.} Write short notes on the following:\pts{$3.5  \times 4 = 14$}

\begin{parts}
  \item Electro-magnetic Radiation.
  \item Spectral Signature.
  \item Air borne and Space borne remote sensing.
  \item Digital Image Processing.
\end{parts}

\medskip\rule{\linewidth}{0.2pt}\medskip




\noindent   \textbf{Question 2.} Explain the concept and scope of Remote Sensing.\pts{14}

\medskip\rule{\linewidth}{0.2pt}\medskip




\noindent   \textbf{Question 3.} Discuss in detail the spectral bands with suitable examples.\pts{14}

\medskip\rule{\linewidth}{0.2pt}\medskip




\noindent   \textbf{Question 4.} What do you mean by Remote Sensing Satellite Platforms and sensors? Elaborate with suitable examples.\pts{14}

\medskip\rule{\linewidth}{0.2pt}\medskip




\noindent   \textbf{Question 5.} Discuss the types and characteristics of Aerial photos.\pts{14}

\medskip\rule{\linewidth}{0.2pt}\medskip




\noindent   \textbf{Question 6.} Differentiate between Visual and Digital Image Processing techniques.\pts{14}

\medskip\rule{\linewidth}{0.2pt}\medskip




\noindent   \textbf{Question 7.} Discuss and analyze various Remote Sensing applications in Resource Mapping and Environmental Monitoring.\pts{14}

\vfill
\rule{\linewidth}{0.4pt}

\begin{center}\small
\href{https://bhu-exam.vercel.app/}{Click here to visit our website}\\[4pt]
\href{https://me-aryan.vercel.app/}{Click here to visit our contributor}
\end{center}

\end{document}
